
\appendix

\section{Units, invariance, and what can be computed}
\label{app:units}

\paragraph{The nat.} Entropy $H=-\sum_i p_i\log p_i$ is defined only once the
base of the logarithm is chosen, and that choice \emph{is} the choice of unit.
With $\log_2$ one gets bits (shannons), with $\log_{10}$ dits, and with the
natural logarithm \emph{nats}. One nat is the information delivered by an
observation that divides the space of possibilities by $e\approx2.718$, where a
bit divides it by two. The conversion is $1~\text{nat}=1/\log 2\approx1.4427$~bit.
This document uses the natural logarithm throughout: entropies, deficits,
divergences, scenario radii and the structural index are in nats.

Three operational readings, from the most useful to the most abstract. First the
exchange rate established in Section~\ref{sec:stress}: a purely directional
scenario of $\eta$ nats is a move of $\sqrt{2\eta}$ sigmas, so $k$~sigmas cost
$k^2/2$ nats. Next, coding: $\KL{p}{q}$ nats is the expected excess code length
when data with law $p$ are encoded with a code optimal for $q$. Finally,
hypothesis testing: a divergence of one nat per observation means that the
likelihood ratio in favour of $p$ grows on average by a factor $e$ with each
observation.

\paragraph{$h$ depends on units, divergences do not.} The differential entropy
of a continuous variable satisfies, for any invertible affine transformation,
\begin{equation}
\label{eq:h-scaling}
h(aX+b) \;=\; h(X) + \log|a| ,
\qquad\text{and in dimension } N,\quad h(AX+b)=h(X)+\log|\det A| .
\end{equation}
Expressing returns in percent rather than decimals therefore adds $N\log 100$ to
$h$: the \emph{level} of $h$, and hence that of $H^{vol}$ and $H^{cov}$, has no
interpretable absolute value. By contrast $\KL{p}{q}$, $I(X;Y)$, $\DD$ and $\JJ$
are invariant under any invertible change of variable, the Jacobians cancelling
between numerator and denominator of the logarithm. \emph{A figure in nats is
comparable across datasets only if it is a divergence.} This is the structural
reason the paper builds a scale-free index.

\paragraph{How is $h(X)$ computed?} In the Gaussian case, in closed form
(Appendix~\ref{app:gauss}). In general, it cannot be without knowing $p$.
Non-parametric estimators exist (\citet{vasicek1976} spacings,
Kozachenko--Leonenko nearest neighbours, kernel methods) but all degrade quickly
beyond three or four dimensions. Estimating a differential entropy in
dimension~20 on a 252-day window is out of reach.

The paper therefore \emph{never} computes $h(p)$. It computes $h(q)$ in closed
form from $\widehat\Sigma$, and the deficit $\DD$ through an approximation of
the \emph{difference} of marginal entropies. A difference of entropies is
dimensionless and estimable; a bare differential entropy is not. The whole
architecture (decompose rather than estimate directly) follows from this
constraint.

\section{Gaussian entropy: proof of equation~\eqref{eq:ceiling}}
\label{app:gauss}

\begin{proof}[Step A: entropy of a multivariate Gaussian]
For $q=\mathcal N(\mu,\Sigma)$ on $\mathbb{R}^N$,
\[
\log q(x) = -\tfrac12\left[N\log 2\pi + \log\det\Sigma + (x-\mu)'\Sigma^{-1}(x-\mu)\right],
\]
so
\[
h(q) = -\mathbb{E}_q[\log q] = \tfrac12 N\log 2\pi + \tfrac12\log\det\Sigma
+ \tfrac12\,\mathbb{E}\!\left[(X-\mu)'\Sigma^{-1}(X-\mu)\right].
\]
The last expectation follows from the trace trick: a scalar equals its trace,
and the trace is cyclic, hence
\[
\mathbb{E}\!\left[(X-\mu)'\Sigma^{-1}(X-\mu)\right]
= \mathbb{E}\!\left[\operatorname{tr}\!\left(\Sigma^{-1}(X-\mu)(X-\mu)'\right)\right]
= \operatorname{tr}(\Sigma^{-1}\Sigma) = \operatorname{tr}(I_N) = N .
\]
Collecting terms, $h(q)=\tfrac12\left[N\log2\pi+\log\det\Sigma+N\right]
=\tfrac12\log\!\left((2\pi e)^N\det\Sigma\right)$.
\end{proof}

\begin{proof}[Step B: factorising the determinant]
With $\Sigma=DRD$, $D=\operatorname{diag}(\sigma_1,\dots,\sigma_N)$ and $R$ the
correlation matrix, multiplicativity of the determinant gives
$\det\Sigma=\det(D)\det(R)\det(D)=\left(\prod_i\sigma_i\right)^2\det R$.
\end{proof}

\begin{proof}[Step C: substitution]
$\tfrac12\log\det\Sigma
=\tfrac12\log\!\left[\left(\prod_i\sigma_i\right)^2\det R\right]
=\sum_i\log\sigma_i+\tfrac12\log\det R$, hence
$h(q)=\tfrac N2\log(2\pi e)+\sum_i\log\sigma_i+\tfrac12\log\det R$, which is
equation~\eqref{eq:ceiling}.
\end{proof}

\paragraph{Sign of the dependence channel.} Hadamard's inequality gives
$\det R\le\prod_i R_{ii}=1$ for a positive semi-definite matrix. A more telling
proof goes through the spectrum: $\det R=\prod_k\lambda_k$ with
$\sum_k\lambda_k=\operatorname{tr}R=N$, and the arithmetic--geometric mean
inequality gives $\prod_k\lambda_k\le\left(\sum_k\lambda_k/N\right)^N=1$, with
equality if and only if every eigenvalue equals~1, that is $R=I$. Hence
$H^{dep}\le0$: \emph{correlation always destroys entropy}.

In the bivariate case, $\det R=1-\rho^2$ and $H^{dep}=\tfrac12\log(1-\rho^2)$;
for $\rho=0.9$ this is $-0.83$~nat. Note that $-H^{dep}$ is then exactly the
mutual information between the two assets: correlation is shared information,
hence destroyed entropy.

\paragraph{Reading note on the tables.} The $H^{cov}$ columns of the v1 results
tables report $H^{vol}+H^{dep}$ and \textbf{omit the constant}
$\tfrac N2\log(2\pi e)$, which equals $28.38$~nats for $N=20$. This constant
offset is harmless since only differences between regimes and changes over time
are interpreted, but it explains why the reported levels are negative.

\section{The entropy deficit: why $\DD=h(q)-h(p)$}
\label{app:deficit}

The Kullback--Leibler divergence between a law $p$ and a reference law $q$ is
$\KL{p}{q}=\int p\log(p/q)$, non-negative and zero if and only if $p=q$ almost
everywhere. It is neither symmetric nor equipped with a triangle inequality: it
is not a distance, hence the word divergence.

\begin{proposition}
\label{prop:deficit}
If $q=\mathcal N(\mu,\Sigma)$ has \emph{the first two moments of $p$}, then
$\KL{p}{q}=h(q)-h(p)$.
\end{proposition}

\begin{proof}
$\KL{p}{q}=-h(p)-\mathbb{E}_p[\log q(X)]$. Now $\log q$ is exactly quadratic in
$x$, so its expectation depends only on the mean and covariance of the law under
which it is evaluated. Since $p$ and $q$ share these two moments,
$\mathbb{E}_p\!\left[(X-\mu)'\Sigma^{-1}(X-\mu)\right]=\operatorname{tr}(\Sigma^{-1}\Sigma)=N
=\mathbb{E}_q[\cdot]$, hence $\mathbb{E}_p[\log q]=\mathbb{E}_q[\log q]=-h(q)$ and
so $\KL{p}{q}=h(q)-h(p)$.
\end{proof}

The moment-matching assumption is essential: without it the identity fails. This
is the whole content of Proposition~\ref{prop:maxent}, read as a measurement
rather than an existence theorem. Economically, $\DD$ is the price in nats of the
lie ``the market is Gaussian'': the apparent uncertainty that the Gaussian
description overstates, because the true law is more concentrated than its
covariance alone suggests.

\section{The structural index is a divergence}
\label{app:index}

\begin{proposition}
\label{prop:index}
With $q_i=\mathcal N(\mu_i,\sigma_i^2)$ the Gaussian marginals with the same
scales as $p$,
\[
\KL{p}{\textstyle\prod_i q_i} \;=\; \DD - H^{dep} \;=\; \JJ .
\]
\end{proposition}

\begin{proof}
By definition,
\[
\begin{aligned}
\KL{p}{\textstyle\prod_i q_i} &= -h(p) - \mathbb{E}_p\!\left[\textstyle\sum_i\log q_i(X_i)\right]
= -h(p) + \sum_i\left[\tfrac12\log(2\pi\sigma_i^2)+\tfrac12\right]\\
&= -h(p) + \sum_i\log\sigma_i + \tfrac N2\log(2\pi e),
\end{aligned}
\]
using $\mathbb{E}_p[(X_i-\mu_i)^2]=\sigma_i^2$, true by construction. Substituting
equation~\eqref{eq:three}, namely
$h(p)=\tfrac N2\log(2\pi e)+\sum_i\log\sigma_i+H^{dep}-\DD$, the scale terms and
the constant cancel and $\DD-H^{dep}$ remains.
\end{proof}

This is an exact identity, not an approximation, and it is what justifies
Definition~\ref{def:J}. Removing the scale channel from \eqref{eq:three} might
look like a fudge; the object that remains turns out to be a proper
Kullback--Leibler divergence, and so inherits three properties for free: it is
non-negative; it is zero if and only if returns are independent \emph{and}
Gaussian, which provides a well-defined rather than arbitrary reference point;
and it is invariant to asset-by-asset changes of units, since $p$ and
$\prod_i q_i$ transform in the same way.

It is useful to set this identity beside the general chain rule
$\KL{p}{\prod_i q_i}=\sum_i\KL{p_i}{q_i}+I(x)$, where
$I(x)=\KL{p}{\prod_i p_i}$ is the multi-information. The first term is the sum of
marginal deficits $\sum_i\Dun(x_i)$, the second the total dependence. The
estimator used in v1 estimates the first term with the \citet{hyvarinen1998}
approximation and replaces the second with its linear part $-H^{dep}$;
non-linear tail dependence is therefore omitted from $\JJ$ and measured
separately through a co-exceedance rate.

\section{Mutual information: the two propositions and an example}
\label{app:mi}

\paragraph{Proof of Proposition~\ref{prop:conditioning}.} By definition of
conditional entropy, $\mathbb{E}_Y[h(X\mid Y)]=h(X,Y)-h(Y)$, and by
\eqref{eq:mi}, $I(X;Y)=h(X)+h(Y)-h(X,Y)$. Adding the two,
$\mathbb{E}_Y[h(X\mid Y)]=h(X)-I(X;Y)$. Non-negativity of $I$ follows from that of
the Kullback--Leibler divergence, of which it is a special case, with equality if
and only if $p_{XY}=p_X\otimes p_Y$.

\paragraph{Proof of Proposition~\ref{prop:constraint}.} Immediate: an extra
constraint can only shrink the feasible set, and the supremum of a functional
over a smaller set cannot increase.

\paragraph{Why the two are not equivalent.} The first quantifies but holds only
on average; the second holds realisation by realisation but concerns the ceiling
and quantifies nothing. A discrete counterexample suffices to show that a
particular realisation can \emph{raise} entropy. Let $X$ be binary with
$\mathbb{P}(X=a)=0.99$, so $H(X)=0.056$~nat. Let a signal $Y$, with probability
$0.98$, remove all doubt ($H(X\mid Y=0)=0$), and with probability $0.02$ lead to
$\mathbb{P}(X=a\mid Y=1)=1/2$, so $H(X\mid Y=1)=\log 2=0.693$~nat, twelve times the
prior entropy. The average is nonetheless
$0.98\times0+0.02\times0.693=0.014$~nat, well below $0.056$, and
$I(X;Y)=0.042$~nat. This is exactly why the paper predicts a jump component that
is \emph{mainly}, not exclusively, one-sided.

\paragraph{Worked Gaussian example.} Let $X\sim\mathcal N(0,\sigma^2)$ be
tomorrow's return and $Y=X+\varepsilon$ a noisy signal,
$\varepsilon\sim\mathcal N(0,\sigma_\varepsilon^2)$ independent. Then
\begin{equation}
\label{eq:mi-gauss}
I(X;Y)=\tfrac12\log\!\left(1+\frac{\sigma^2}{\sigma_\varepsilon^2}\right),
\qquad
\operatorname{Var}(X\mid Y)=\frac{\sigma^2\sigma_\varepsilon^2}{\sigma^2+\sigma_\varepsilon^2},
\end{equation}
the posterior variance not depending on the realisation $y$. With
$\sigma=1\%$ a day:

\begin{center}
\small
\begin{tabular}{@{}lrrrr@{}}
\toprule
Signal quality & $\sigma_\varepsilon$ & $I(X;Y)$ (nats) & posterior vol. & $\sqrt{2\eta}$ equivalent \\
\midrule
sharp    & $0.5\%$ & $0.805$ & $0.45\%$ & $1.27\sigma$ \\
moderate & $1.0\%$ & $0.347$ & $0.71\%$ & $0.83\sigma$ \\
vague    & $3.0\%$ & $0.053$ & $0.95\%$ & $0.32\sigma$ \\
\bottomrule
\end{tabular}
\end{center}

A factor of fifteen separates the size of the entropy jump caused by decisive
news from that caused by a rumour. The last column translates each jump onto the
severity scale of Section~\ref{sec:stress}: it is the directional move a
scenario of the same radius would produce.
